\documentclass[aspectratio=169,11pt]{beamer}
% Shared CMPSC 480 Beamer preamble.
\usepackage[T1]{fontenc}
\usepackage[utf8]{inputenc}
\usepackage{lmodern}
\usepackage{amsmath,amssymb,mathtools}
\usepackage{booktabs,array,tabularx}
\usepackage{xcolor}
\usepackage{listings}
\usepackage{tikz}
\usetikzlibrary{arrows.meta,positioning,shapes.geometric}

% Ignore only negligible TeX box diagnostics that are below visual tolerance.
% Larger layout defects still appear in the build log and in rendered QA.
\vfuzz=3pt
\hbadness=2000

\definecolor{courseink}{HTML}{111111}
\definecolor{coursegray}{HTML}{EDEDED}
\definecolor{courserule}{HTML}{B8BCC4}
\definecolor{courseblue}{HTML}{3D8DFF}
\definecolor{courselightblue}{HTML}{D0EDFA}
\definecolor{coursemuted}{HTML}{5C6470}
\definecolor{coursegreen}{HTML}{0B7A53}
\definecolor{coursered}{HTML}{B3261E}

\usetheme{default}
\usefonttheme{professionalfonts}
\setbeamertemplate{navigation symbols}{}
\setbeamersize{text margin left=0.65cm,text margin right=0.65cm}

\setbeamercolor{normal text}{fg=courseink,bg=white}
\setbeamercolor{frametitle}{fg=courseink,bg=white}
\setbeamercolor{title}{fg=courseink,bg=white}
\setbeamercolor{subtitle}{fg=coursemuted,bg=white}
\setbeamercolor{structure}{fg=courseblue}
\setbeamercolor{alerted text}{fg=coursered}
\setbeamercolor{block title}{fg=courseink,bg=courselightblue}
\setbeamercolor{block body}{fg=courseink,bg=coursegray}
\setbeamercolor{example text}{fg=coursegreen}

\setbeamerfont{title}{size=\fontsize{25}{29}\selectfont,series=\bfseries}
\setbeamerfont{subtitle}{size=\fontsize{13}{16}\selectfont}
\setbeamerfont{frametitle}{size=\fontsize{19}{22}\selectfont,series=\bfseries}
\setbeamerfont{normal text}{size=\fontsize{11}{14}\selectfont}
\setbeamerfont{footline}{size=\fontsize{7.5}{9}\selectfont}

\setbeamertemplate{frametitle}{%
  \nointerlineskip
  % Use content-driven height so a long assessment title can wrap without
  % being clipped by a fixed-height title bar.
  \begin{beamercolorbox}[wd=\paperwidth,sep=0.17cm,leftskip=0.48cm,rightskip=0.48cm]{frametitle}
    \insertframetitle
  \end{beamercolorbox}
  \vspace{-0.08cm}
  {\color{courserule}\rule{\textwidth}{0.35pt}}
}

\setbeamertemplate{footline}{%
  \leavevmode
  \begin{beamercolorbox}[wd=\paperwidth,ht=2.6ex,dp=1.1ex,leftskip=.65cm,rightskip=.65cm]{author in head/foot}
    \color{coursemuted}\insertshorttitle\hfill\insertframenumber/\inserttotalframenumber
  \end{beamercolorbox}
}

\setbeamertemplate{itemize item}{\color{courseblue}\small$\blacktriangleright$}
\setbeamertemplate{itemize subitem}{\color{courseblue}\scriptsize$\bullet$}
\setbeamertemplate{enumerate item}{\color{courseblue}\insertenumlabel.}

\lstdefinestyle{coursepython}{
  language=Python,
  basicstyle=\ttfamily\fontsize{8.5}{10}\selectfont,
  keywordstyle=\color{courseblue}\bfseries,
  commentstyle=\color{coursemuted}\itshape,
  stringstyle=\color{coursegreen},
  showstringspaces=false,
  columns=fullflexible,
  keepspaces=true,
  frame=single,
  rulecolor=\color{courserule},
  backgroundcolor=\color{coursegray},
  breaklines=true,
  tabsize=4,
  xleftmargin=0.15cm,
  xrightmargin=0.15cm,
  framexleftmargin=0.15cm,
  framexrightmargin=0.15cm
}
\lstset{style=coursepython}

\newcommand{\points}[1]{\hfill{\color{courseblue}\textbf{[#1 points]}}}
\newcommand{\deliverable}[1]{\textbf{\color{courseblue}#1}}
\newcommand{\code}[1]{\texttt{#1}}
\newcommand{\keyidea}[1]{%
  \begin{beamercolorbox}[rounded=false,sep=0.55em,wd=\linewidth]{block body}
    \textbf{Key idea.} #1
  \end{beamercolorbox}
}
\newcommand{\answerlabel}{\textbf{\color{coursegreen}Answer.}\ }
\newcommand{\gradinglabel}{\textbf{\color{courseblue}Grading guidance.}\ }

\newenvironment{questionblock}[1]
  {\begin{block}{#1}}
  {\end{block}}

\newenvironment{solutionblock}[1]
  {\begin{block}{#1}}
  {\end{block}}

\AtBeginSection[]{
  \begin{frame}
    \vfill
    {\usebeamerfont{title}\color{courseink}\insertsectionhead\par}
    \vspace{0.4cm}
    {\color{courseblue}\rule{0.32\paperwidth}{3pt}}
    \vfill
  \end{frame}
}


% CMPSC480_TOTAL_POINTS: 10
% CMPSC480_ITEM_IDS: 1,2,3,4
% CMPSC480_SUBITEM_POINTS: 1=2,2=5,3=1,4=2
\title[CMPSC 480 - Week 1 Solutions]{Week 1 Code Lab - Complete Solutions}
\subtitle{Environment Setup and Your First Agent/Problem Skeleton}
\author{CMPSC 480 - Instructor Solution Deck}
\date{}

\begin{document}

\begin{frame}
  \titlepage
  \vspace{-0.35cm}
  \begin{center}
    \color{coursered}\bfseries Instructor material - do not distribute with the question deck
  \end{center}
\end{frame}

\begin{frame}{Solution overview}
  \answerlabel The reference implementation uses:
  \begin{itemize}
    \item a tuple \((row, column)\) as the immutable state;
    \item an \texttt{Enum} for the four actions;
    \item a validated \texttt{GridProblem} with deterministic action order;
    \item an abstract generic \texttt{Agent} contract;
    \item \texttt{unittest} for normal, boundary, illegal, and malformed cases; and
    \item a verification script that fails with a nonzero process status.
  \end{itemize}

  \vspace{0.25cm}
  Equivalent designs earn full credit when they satisfy the same observable
  contract.
\end{frame}

\begin{frame}{Reference repository structure}
  \begin{block}{Complete submission}
    \ttfamily
    week01-agent-skeleton/\par
    \hspace{0.45cm}README.md\par
    \hspace{0.45cm}verify\_environment.py\par
    \hspace{0.45cm}src/\par
    \hspace{0.9cm}\_\_init\_\_.py\par
    \hspace{0.9cm}agent.py\par
    \hspace{0.9cm}problem.py\par
    \hspace{0.45cm}tests/\par
    \hspace{0.9cm}test\_problem.py
  \end{block}
  \gradinglabel Accept a \texttt{pyproject.toml} or package-based structure if
  imports and the documented test command work from a clean checkout.
\end{frame}

\begin{frame}[fragile]{Solution 1 - environment verification imports}
\begin{lstlisting}
from __future__ import annotations

import subprocess
import sys
from importlib.metadata import PackageNotFoundError, version


def package_version(name: str) -> str:
    try:
        return version(name)
    except PackageNotFoundError as exc:
        raise RuntimeError(
            f"required package is not installed: {name}"
        ) from exc
\end{lstlisting}
\end{frame}

\begin{frame}[fragile]{Solution 1 - Git command helper}
\begin{lstlisting}
def git_output(*args: str) -> str:
    completed = subprocess.run(
        ["git", *args],
        check=True,
        capture_output=True,
        text=True,
    )
    return completed.stdout.strip()
\end{lstlisting}
\answerlabel \texttt{check=True} converts a failed Git command into a nonzero
program outcome instead of silently accepting empty output.
\end{frame}

\begin{frame}[fragile]{Solution 1 - version and package checks}
\begin{lstlisting}
def main() -> None:
    if sys.version_info < (3, 10):
        raise RuntimeError("Python 3.10 or newer is required")

    print(f"PYTHON={sys.version.split()[0]}")
    print(f"NUMPY={package_version('numpy')}")
    print(f"JUPYTER={package_version('jupyterlab')}")
\end{lstlisting}
\answerlabel Stable labels make the output readable by both a person and a
simple grading script.
\end{frame}

\begin{frame}[fragile]{Solution 1 - Git checks and entry point}
\begin{lstlisting}
    # Continue inside main().
    if git_output("rev-parse", "--is-inside-work-tree") != "true":
        raise RuntimeError("current directory is not a Git work tree")

    commits = int(git_output("rev-list", "--count", "HEAD"))
    if commits < 1:
        raise RuntimeError("the repository has no commits")
    print(f"GIT_COMMITS={commits}")


if __name__ == "__main__":
    main()
\end{lstlisting}
\end{frame}

\begin{frame}{Solution 1 - Expected behavior (2 points)}
  \begin{block}{Representative successful output}
    \ttfamily
    PYTHON=3.11.0\par
    NUMPY=2.1.0\par
    JUPYTER=4.2.0\par
    GIT\_COMMITS=3
  \end{block}
  \begin{itemize}
    \item Exact installed versions may differ.
    \item Missing software, an old Python version, a nonrepository directory, or
      zero commits must produce a nonzero exit status.
  \end{itemize}
  \gradinglabel Award full credit when the output is stable and every required
  condition affects process success.
\end{frame}

\begin{frame}[fragile]{Solution 2 - immutable state and action model}
\begin{lstlisting}
from __future__ import annotations

from enum import Enum
from typing import TypeAlias

State: TypeAlias = tuple[int, int]


class Action(Enum):
    UP = (-1, 0)
    DOWN = (1, 0)
    LEFT = (0, -1)
    RIGHT = (0, 1)
\end{lstlisting}
  \answerlabel A two-integer tuple has value equality, stable hashing, and no
  mutating methods. It can therefore serve as a set member or dictionary key.
\end{frame}

\begin{frame}[fragile]{Solution 2 - class and constructor signature}
\begin{lstlisting}
class GridProblem:
    ACTION_ORDER = (
        Action.UP, Action.DOWN, Action.LEFT, Action.RIGHT
    )

    def __init__(
        self,
        rows: int,
        columns: int,
        start: State,
        goal: State,
        blocked: frozenset[State] = frozenset(),
    ) -> None:
\end{lstlisting}
\answerlabel A fixed action order makes successor generation deterministic,
which simplifies testing and debugging.
\end{frame}

\begin{frame}[fragile]{Solution 2 - constructor state}
\begin{lstlisting}
        # Continue inside GridProblem.__init__.
        if rows < 1 or columns < 1:
            raise ValueError("grid dimensions must be positive")
        self.rows = rows
        self.columns = columns
        self.start = start
        self.goal = goal
        self.blocked = frozenset(blocked)
\end{lstlisting}
\answerlabel Copying \texttt{blocked} into a \texttt{frozenset} prevents later
caller mutation from changing the problem definition.
\end{frame}

\begin{frame}[fragile]{Solution 2 - constructor domain validation}
\begin{lstlisting}
        # Continue inside GridProblem.__init__.
        self._validate_state(start)
        self._validate_state(goal)
        for state in self.blocked:
            self._validate_coordinate(state)
        if start in self.blocked or goal in self.blocked:
            raise ValueError("start and goal must be traversable")
\end{lstlisting}
\answerlabel Every configured coordinate is checked once at the public boundary.
\end{frame}

\begin{frame}[fragile]{Solution 2 - state-validation helpers}
\begin{lstlisting}
    def _validate_coordinate(self, state: State) -> None:
        if (
            not isinstance(state, tuple)
            or len(state) != 2
            or any(type(value) is not int for value in state)
        ):
            raise TypeError("state must be a tuple of two integers")
        row, column = state
        if not (0 <= row < self.rows and 0 <= column < self.columns):
            raise ValueError(f"state is outside the grid: {state}")

    def _validate_state(self, state: State) -> None:
        self._validate_coordinate(state)
        if state in self.blocked:
            raise ValueError(f"state is blocked: {state}")

    def initial_state(self) -> State:
        return self.start
\end{lstlisting}
\end{frame}

\begin{frame}[fragile]{Solution 2 - legal actions}
\begin{lstlisting}
    def _destination(self, state: State, action: Action) -> State:
        delta_row, delta_column = action.value
        return state[0] + delta_row, state[1] + delta_column

    def actions(self, state: State) -> list[Action]:
        self._validate_state(state)
        legal: list[Action] = []
        for action in self.ACTION_ORDER:
            candidate = self._destination(state, action)
            try:
                self._validate_state(candidate)
            except ValueError:
                continue
            legal.append(action)
        return legal
\end{lstlisting}
\answerlabel Iterating through \texttt{ACTION\_ORDER} makes tie behavior
deterministic and testable.
\end{frame}

\begin{frame}[fragile]{Solution 2 - validated transition}
\begin{lstlisting}
    def result(self, state: State, action: Action) -> State:
        self._validate_state(state)
        if not isinstance(action, Action):
            raise TypeError("action must be an Action value")
        if action not in self.actions(state):
            raise ValueError(f"illegal action {action.name} at {state}")
        return self._destination(state, action)
\end{lstlisting}
\answerlabel The returned tuple is a new value. The caller's state remains
unchanged and safe as an explored-set key.
\end{frame}

\begin{frame}[fragile]{Solution 2 - validated goal test}
\begin{lstlisting}
    def goal_test(self, state: State) -> bool:
        self._validate_state(state)
        return state == self.goal
\end{lstlisting}
\answerlabel Validation prevents a malformed coordinate from silently becoming a
normal negative result.
\end{frame}

\begin{frame}[fragile]{Solution 2 - verified step cost}
\begin{lstlisting}
    def step_cost(
        self,
        state: State,
        action: Action,
        next_state: State,
    ) -> float:
        expected = self.result(state, action)
        self._validate_state(next_state)
        if next_state != expected:
            raise ValueError(
                f"{next_state} is not the result of {action.name} at {state}"
            )
        return 1.0
\end{lstlisting}
  \answerlabel Reusing \texttt{result} centralizes the legal-transition contract
  and prevents cost calculation for an impossible edge.
\end{frame}

\begin{frame}[fragile]{Solution 3 - Agent contract (1 point)}
\begin{lstlisting}
from abc import ABC, abstractmethod
from typing import Generic, TypeVar

PerceptT = TypeVar("PerceptT")
ActionT = TypeVar("ActionT")


class Agent(ABC, Generic[PerceptT, ActionT]):
    """Base interface for agents that map percepts to actions."""

    @abstractmethod
    def act(self, percept: PerceptT) -> ActionT:
        """Return the next action for the supplied percept."""
        raise NotImplementedError
\end{lstlisting}
  \answerlabel The abstract method makes incomplete use fail immediately while
  preserving a stable Week 2 extension point.
\end{frame}

\begin{frame}[fragile]{Solution 4 - Contract tests and repository quality (2 points)}
\begin{lstlisting}
import unittest

from src.problem import Action, GridProblem


class GridProblemTests(unittest.TestCase):
    def setUp(self) -> None:
        self.problem = GridProblem(
            rows=3,
            columns=3,
            start=(0, 0),
            goal=(2, 2),
            blocked=frozenset({(1, 1)}),
        )
\end{lstlisting}
\answerlabel A shared fixture keeps every test focused on one behavior.
\end{frame}

\begin{frame}[fragile]{Solution 4 - normal and boundary tests}
\begin{lstlisting}
    def test_corner_actions_are_deterministic(self) -> None:
        self.assertEqual(
            self.problem.actions((0, 0)),
            [Action.DOWN, Action.RIGHT],
        )

    def test_known_transition(self) -> None:
        self.assertEqual(
            self.problem.result((0, 0), Action.RIGHT),
            (0, 1),
        )
\end{lstlisting}
\gradinglabel Award one point for focused behavioral tests and one point for
reproducible documentation and repository quality.
\end{frame}

\begin{frame}[fragile]{Solution 4 - illegal, goal, and hashing tests}
\begin{lstlisting}
    def test_illegal_action_is_rejected(self) -> None:
        with self.assertRaises(ValueError):
            self.problem.result((0, 0), Action.UP)

    def test_goal_test_and_hashability(self) -> None:
        self.assertFalse(self.problem.goal_test((0, 0)))
        self.assertTrue(self.problem.goal_test((2, 2)))
        explored = {(0, 0), (0, 1)}
        self.assertIn((0, 1), explored)
\end{lstlisting}
\end{frame}

\begin{frame}[fragile]{Solution 4 - degenerate and malformed cases}
\begin{lstlisting}
    def test_one_cell_grid_has_no_actions(self) -> None:
        problem = GridProblem(1, 1, (0, 0), (0, 0))
        self.assertEqual(problem.actions((0, 0)), [])

    def test_malformed_state_is_rejected(self) -> None:
        with self.assertRaises(ValueError):
            self.problem.goal_test((3, 0))


if __name__ == "__main__":
    unittest.main()
\end{lstlisting}
\end{frame}

\begin{frame}{Edge-case reasoning}
  \begin{columns}[T,onlytextwidth]
    \begin{column}{0.47\textwidth}
      \begin{block}{One-cell grid}
        The state may be both initial and goal. The legal-action list is empty,
        which is distinct from a failure value.
      \end{block}
      \begin{block}{Blocked cells}
        Configuration may name blocked coordinates, but no public operation may
        accept one as a traversable state.
      \end{block}
    \end{column}
    \begin{column}{0.47\textwidth}
      \begin{block}{Illegal action}
        Returning the unchanged state would hide a caller defect. The reference
        contract raises \texttt{ValueError}.
      \end{block}
      \begin{block}{Stable hashing}
        Every transition creates a new tuple. Existing state keys remain valid
        in sets and dictionaries.
      \end{block}
    \end{column}
  \end{columns}
\end{frame}

\begin{frame}{Rubric solution - Correctness (8 points)}
  \scriptsize
  \begin{tabularx}{\textwidth}{>{\bfseries}p{0.25\textwidth} c X}
    \toprule
    Category & Points & Full-credit evidence \\
    \midrule
    Environment verification & 2 & All versions, Git work-tree status, positive
      commit count, and nonzero failure status \\
    Immutable state model & 1 & Stable equality and hashing; no mutation \\
    Legal actions and transitions & 2 & Correct deterministic behavior at
      boundaries and blocked cells \\
    Goal test and step cost & 1 & Boolean goal result; positive cost only for the
      matching legal transition \\
    Defensive validation & 1 & Precise type/domain/action rejection \\
    Agent abstraction & 1 & Abstract, documented, type-parameterized contract \\
    \bottomrule
  \end{tabularx}
\end{frame}

\begin{frame}{Rubric solution - Tests and quality (2 points)}
  \scriptsize
  \begin{tabularx}{\textwidth}{>{\bfseries}p{0.25\textwidth} c X}
    \toprule
    Category & Points & Full-credit evidence \\
    \midrule
    Focused tests & 1 & Normal transition, boundary actions, illegal action,
      goal cases, hashability, no-action state, and malformed state \\
    Documentation and repository & 1 & Public docstrings, complete type hints,
      reproducible README, clean tree, and meaningful commit \\
    \midrule
    \textbf{Total} & \textbf{10} & \\
    \bottomrule
  \end{tabularx}

  \vspace{0.35cm}
  \gradinglabel Deduct only once for one underlying defect that causes several
  downstream test failures.
\end{frame}

\begin{frame}{Reflection exemplar}
  \small
  \answerlabel The initial-state component maps to the immutable tuple returned
  by \texttt{initial\_state}. The action function maps to
  \texttt{actions(state)}, whose deterministic list contains exactly the moves
  that stay in bounds and avoid blocked cells. The transition model maps to
  \texttt{result}, which validates the current state and action before creating
  a new tuple. The goal test compares a validated state with the configured
  goal, and the path-cost model assigns 1.0 to each verified legal transition.

  \vspace{0.3cm}
  One ambiguity was whether an illegal move should leave the state unchanged.
  I chose an exception because the method contract accepts only legal actions;
  silently staying in place would hide a caller defect. A tuple state has value
  equality and stable hashing, so Week 2 can store it safely in an explored set.
\end{frame}

\begin{frame}{Instructor verification commands}
  Run from a clean checkout:

  \vspace{0.35cm}
  \begin{block}{Environment}
    \ttfamily python verify\_environment.py
  \end{block}
  \begin{block}{Tests}
    \ttfamily python -m unittest discover -s tests -v
  \end{block}
  \begin{block}{Repository evidence}
    \ttfamily git status --short\par
    git log -1 --oneline
  \end{block}

  \keyidea{The reference implementation contains no search algorithm; Week 1
  establishes the abstraction that Week 2 will extend.}
\end{frame}

\end{document}
